\documentclass[12pt]{amsart}
% Default margins are too wide all the way around. I reset them here
\setlength{\hoffset}{-.75in}
\setlength{\textwidth}{6.5in}
\setlength{\voffset}{-.5in}
\setlength{\textheight}{9.0in}
\usepackage{amsmath}
\usepackage{amssymb}
\usepackage{amsthm}
\usepackage{pgfplots}
\usepackage{enumerate}
\usepackage{tikz}
\usetikzlibrary{shadows,arrows,arrows.meta,shapes,positioning,calc}% Required for the Circle and Triangle arrow tips
\usepackage{setspace}
\usepackage{siunitx}
\usepackage{enumitem}
\usepackage{tcolorbox}
\usepackage{array}
\usepackage{microtype}
\usepackage{hyperref}
\usepackage{upgreek}

\theoremstyle{conjecture}
\newtheorem{conjecture}{Conjecture}

\theoremstyle{definition}
\newtheorem{definition}{Definition}

\theoremstyle{question}
\newtheorem{question}{Question}

\theoremstyle{theorem}
\newtheorem{theorem}{Theorem}

\newenvironment{solution}
{\begin{proof}[Solution]}
{\end{proof}}

\newcommand{\sectionrule}{%
    \vspace{0.3em}
    \hrule
    \vspace{0.8em}
}

\newcommand{\vct}[1]{\underline{#1}}
\newcommand{\sys}[2]{\begin{aligned}#1\\#2\end{aligned}}
\newcommand{\rulehead}[1]{\vspace{2pt}\hrule\vspace{5pt}{\large\bfseries #1}\par}

\newcommand{\phaseplot}[1]{%
    \par\medskip
    \noindent\begin{minipage}{\linewidth}
    \centering
    \begin{tikzpicture}
    \begin{axis}[
    width=0.61\linewidth,height=4.75cm,
    axis lines=middle,
    xmin=-3.3,xmax=3.3,ymin=-1.12,ymax=1.16,
    xtick={-3,-2,-1,1,2,3},ytick={-1,-0.5,0.5,1},
    minor tick num=4,
    xlabel={$x$},ylabel={},
    title={\small functions},title style={yshift=-0.6em},
    tick label style={font=\small},
    every axis x label/.style={at={(ticklabel* cs:1)},anchor=west},
    axis line style={gray!70},tick style={gray!70},
    clip=false]
    \addplot[blue!55!gray,thick,domain=-3.3:3.3,samples=150]{sin(deg(x))};
    \addplot[orange!85!yellow!70!black,thick,domain=-3.2:3.2]{#1};
    \end{axis}
    \end{tikzpicture}
    \end{minipage}
    \par\medskip}

\title{Dynamical Systems:\\Class September 15, 2026\\ 6. 2D Nonlinear Systems and Fixed Points}
\author{Rob Tompkins}
\date{\today}


\begin{document}

    \maketitle
    {\large\bfseries Teams}

    Post screenshots of your work to the course Google Drive today: Link to drive folder. Include words, labels, and other short notes on your whiteboard that might make those solutions useful to you or your classmates. Name your images using the format:
    \begin{itemize}\item \texttt{teamXX-YY.EXT}\end{itemize}
    where XX is your two-digit team number (e.g. 06 if you are team 6), YY is the two digit number of image for today (e.g. 01 for the first image, 02 for the second, and so on.), and EXT is the extension for the image filetype you used.
    \qquad\\
    \qquad\\
    {\large\bfseries Team Activity}

    \begin{enumerate}[leftmargin=2em]
      \item[(1.)] (6.4.2) Consider the system
      \[
          \dot{x}=x(3-2x-y),
          \qquad
          \dot{y}=y(2-x-y),
          \qquad x,y\geq 0.
      \]
      \begin{enumerate}[label=(\alph*),leftmargin=2.2em]
        \item Find the fixed points.
        \item Draw the nullclines on the $xy$-plane.
        \item Add a representative vector in each region of phase space.
        \item Compute the Jacobian matrix.
        \item Classify the fixed points. Compare your classification to the behavior
        of the vector field near each fixed point in your diagram from part (c).
        \item Sketch a plausible phase portrait for the system. Compare your sketch
        to a numerically computed phase portrait.
      \end{enumerate}

      \begin{solution}
          Write
          \[
              f(x,y)=x(3-2x-y),\qquad g(x,y)=y(2-x-y).
          \]

          \begin{enumerate}[label=(\alph*),leftmargin=2.2em]
        \item The fixed points are the intersections of
        \[
            x=0\quad\hbox{or}\quad 3-2x-y=0
        \]
        with
        \[
            y=0\quad\hbox{or}\quad 2-x-y=0.
        \]
        Thus
        \[
            \boxed{(0,0),\quad(0,2),\quad
            \left(\frac32,0\right),\quad(1,1)}.
        \]
        For the interior point, subtracting
        $2x+y=3$ and $x+y=2$ gives $x=1$, hence $y=1$.

        \item The $x$-nullclines are
        \[
            x=0,\qquad y=3-2x,
        \]
        and the $y$-nullclines are
        \[
            y=0,\qquad y=2-x.
        \]
        In the first quadrant the oblique $x$-nullcline runs from $(0,3)$
        to $(3/2,0)$, while the oblique $y$-nullcline runs from $(0,2)$ to
        $(2,0)$.  They cross at $(1,1)$.

        \item Away from the axes, the signs are determined by
        \[
            \begin{array}{c|c}
                \text{condition}&\text{sign}\rule{0pt}{2.4ex}\\ \hline
                y<3-2x&\dot x>0\\
                y>3-2x&\dot x<0\\
                y<2-x&\dot y>0\\
                y>2-x&\dot y<0
            \end{array}
        \]
        Therefore the four possible directions are
        $\nearrow$ for $(+,+)$, $\searrow$ for $(+,-)$,
        $\nwarrow$ for $(-,+)$, and $\swarrow$ for $(-,-)$.
        For example, representatives are
        \[
            \begin{array}{c|c|c}
                (x,y)&(\operatorname{sgn}\dot x,\operatorname{sgn}\dot y)&
                \text{direction}\\ \hline
                (1/2,1/2)&(+,+)&\nearrow\\
                (1/2,7/4)&(+,-)&\searrow\\
                (5/4,3/5)&(-,+)&\nwarrow\\
                (1/2,5/2)&(-,-)&\swarrow
            \end{array}
        \]

        \item Differentiating $f$ and $g$ gives
        \[
            \boxed{J(x,y)=
                \begin{bmatrix}
                    3-4x-y&-x\\
                    -y&2-x-2y
                \end{bmatrix}}.
        \]

        \item Evaluation at the fixed points gives
        \[
            \begin{array}{c|c|c|c}
                (x,y)&J(x,y)&\text{eigenvalues}&\text{classification}\\ \hline
            (0,0)&\begin{bmatrix}3&0\\0&2\end{bmatrix}
                &3,2&\text{unstable node}\\[1.2ex]
                (0,2)&\begin{bmatrix}1&0\\-2&-2\end{bmatrix}
                &1,-2&\text{saddle}\\[1.2ex]
                (3/2,0)&\begin{bmatrix}-3&-3/2\\0&1/2\end{bmatrix}
                &-3,1/2&\text{saddle}\\[1.2ex]
                (1,1)&\begin{bmatrix}-2&-1\\-1&-1\end{bmatrix}
                &\dfrac{-3\pm\sqrt5}{2}&\text{stable node}
            \end{array}
        \]
        These classifications agree with the local arrow directions from
        part (c): arrows leave $(0,0)$, enter and leave each saddle along
        different directions, and point toward $(1,1)$ from all nearby
        sectors.

        \item The coordinate axes are invariant.  On the positive $x$-axis,
        solutions approach $(3/2,0)$; on the positive $y$-axis, solutions
        approach $(0,2)$.  Interior trajectories cannot cross either axis
        and are directed toward the stable coexistence equilibrium $(1,1)$.
        A numerical phase portrait should therefore show the two boundary
        saddles, their invariant-axis trajectories, and interior trajectories
        curving into $(1,1)$.
          \end{enumerate}
      \end{solution}

      \item[(2.)] \textbf{Analyzing a 2D System:} Consider a 2D red fox--coyote
      system (let ``1'' denote red foxes and ``2'' denote coyotes):
      \begin{align*}
          \frac{dN_1}{dt}
          &=r_1N_1\left(1-\frac{N_1}{K_1}\right)-\alpha_1N_1N_2,\\
          \frac{dN_2}{dt}
          &=r_2N_2\left(1-\frac{N_2}{K_2}\right)-\alpha_2N_1N_2.
      \end{align*}

      \begin{enumerate}[label=(\alph*),leftmargin=2.2em]
        \item Using $N_1=A_1x$, $N_2=A_2y$, and $t=T\tau$, show that the
        system can be nondimensionalized to give
        \begin{align*}
            \frac{dx}{d\tau}&=x(1-x)-\beta_1xy,\\
            \frac{dy}{d\tau}&=\rho y(1-y)-\beta_2xy,
        \end{align*}
        where
        \[
            \rho=\frac{r_2}{r_1},
            \qquad
            \beta_1=\frac{\alpha_1K_2}{r_1},
            \qquad
            \beta_2=\frac{\alpha_2K_1}{r_1}.
        \]
        \begin{enumerate}[label=\roman*.,leftmargin=2em]
            \item Show that the nullclines of this nondimensional system are
            \begin{align*}
                x\text{ nullclines:}&\quad x=0,
                &y&=\frac{1}{\beta_1}(1-x),\\
                y\text{ nullclines:}&\quad y=0,
                &y&=1-\frac{\beta_2}{\rho}x.
            \end{align*}
            \item There are three equilibria whose locations are unaffected by the
            parameters. What are they? Interpret these equilibria in terms of
            foxes and coyotes.
            \item Find conditions to ensure that there is exactly one crossing of
            the nullclines in the first quadrant (not on the axes). How would you
            interpret these conditions in terms of foxes and coyotes?
        \end{enumerate}

        \item Show that the Jacobian for the dimensionless system is
        \[
            J=
            \begin{bmatrix}
                1-2x-\beta_1y & -\beta_1x\\
                -\beta_2y & \rho(1-2y)-\beta_2x
            \end{bmatrix}.
        \]
        \begin{enumerate}[label=\roman*.,leftmargin=2em]
            \item Determine the linear stability of equilibrium $e_1=(1,0)$ and
            sketch the phase portraits nearby. Are there any cases where the
            linear stability cannot tell us about the nonlinear stability of this
            equilibrium? Relate your analysis back to foxes and coyotes.
            \item Determine the linear stability of equilibrium $e_2=(0,1)$ and
            sketch the phase portraits nearby. Are there any cases where the
            linear stability cannot tell us about the nonlinear stability of this
            equilibrium? Relate your analysis back to foxes and coyotes.
            \item Determining the linear stability of equilibrium $e_4$ (the
            crossing point) is messy without a computer. Instead, use the
            nullclines around this point to determine the stability when
            $\beta_1<1$ and $\beta_2<\rho$, and when $\beta_1>1$ and
            $\beta_2>\rho$.
        \end{enumerate}

        \item Conceptual questions about the model:
        \begin{enumerate}[label=\roman*.,leftmargin=2em]
            \item What are some assumptions that are explicit or implicit in this
            model?
            \item Briefly describe each of the terms in the model. Which terms give
            the intraspecies competition? Which give the interspecies
            competition?
            \item Which parameters are likely to be hard to measure/constrain?
            Which would be easy to measure/constrain?
        \end{enumerate}

        \item Conceptual questions about results/analysis:
            \begin{enumerate}[label=\roman*.,leftmargin=2em]
                \item In New York's Adirondack region, coyotes prey on red foxes and the
                NYS Department of Environmental Conservation reports that red foxes
                ``avoid coyote territories completely or reside on the periphery of
                established coyote territories.'' Is there any inconsistency with
                this fact and the model we just analyzed? If no, explain why. If yes,
                can you adjust the model to fix it?
                \item What in the above analysis would you change if we replaced the red
                foxes with bobcats? The Adirondack Ecological Center site on bobcats
                will be helpful:
                \url{https://www.esf.edu/aec/adks/mammals/bobcat.htm}.
                \item What in the above analysis would you change if we replaced the
                coyotes with gray foxes? The Adirondack Ecological Center site on gray
                foxes will be helpful:
                \url{https://www.esf.edu/aec/adks/mammals/gray_fox.htm}.
            \end{enumerate}
        \end{enumerate}

      \begin{solution}
          Put
          \[
              F(x,y)=x(1-x)-\beta_1xy,
              \qquad
              G(x,y)=\rho y(1-y)-\beta_2xy.
          \]

          \begin{enumerate}[label=(\alph*),leftmargin=2.2em]
          \item Since $N_1=A_1x$, $N_2=A_2y$, and $t=T\tau$,
          \[
              \frac{dN_1}{dt}=\frac{A_1}{T}\frac{dx}{d\tau},
              \qquad
              \frac{dN_2}{dt}=\frac{A_2}{T}\frac{dy}{d\tau}.
          \]
          Substitution into the dimensional equations and division by $A_1/T$
          and $A_2/T$, respectively, yields
          \begin{align*}
              \frac{dx}{d\tau}
              &=Tr_1x\left(1-\frac{A_1}{K_1}x\right)
              -T\alpha_1A_2xy,\\
              \frac{dy}{d\tau}
              &=Tr_2y\left(1-\frac{A_2}{K_2}y\right)
              -T\alpha_2A_1xy.
          \end{align*}
          Choose
          \[
              A_1=K_1,\qquad A_2=K_2,
              \qquad T=\frac1{r_1}.
          \]
          Then
          \[
              \frac{dx}{d\tau}=x(1-x)-\beta_1xy,
              \qquad
              \frac{dy}{d\tau}=\rho y(1-y)-\beta_2xy,
          \]
          with the stated definitions of $\rho$, $\beta_1$, and $\beta_2$.

          \begin{enumerate}[label=\roman*.,leftmargin=2em]
            \item Factoring gives
            \[
                F=x(1-x-\beta_1y),
                \qquad
                G=y\bigl(\rho(1-y)-\beta_2x\bigr).
            \]
            Hence
            \[
                \boxed{x=0\ \text{or}\ y=\frac{1-x}{\beta_1}},
                \qquad
                \boxed{y=0\ \text{or}\ y=1-\frac{\beta_2}{\rho}x}.
            \]

            \item The parameter-independent equilibria are
            \[
                e_0=(0,0),\qquad e_1=(1,0),\qquad e_2=(0,1).
            \]
            They represent extinction of both species, foxes alone at their
            carrying capacity $N_1=K_1$, and coyotes alone at their carrying
            capacity $N_2=K_2$, respectively.

            \item Solving the two oblique-nullcline equations gives
            \[
                e_4=(x_*,y_*),\qquad
                x_*=\frac{\rho(1-\beta_1)}{\rho-\beta_1\beta_2},
                \qquad
                y_*=\frac{\rho-\beta_2}{\rho-\beta_1\beta_2}.
            \]
            This point is strictly in the first quadrant precisely when
            \[
                \boxed{\beta_1<1\ \text{and}\ \beta_2<\rho}
                \quad\text{or}\quad
                \boxed{\beta_1>1\ \text{and}\ \beta_2>\rho}.
            \]
            In the first regime, each species limits itself more strongly than
            it limits the other species, so stable coexistence is possible.  In
            the second, interspecific effects are relatively stronger; the
            crossing still exists, but it separates the basins in a competitive
            exclusion regime.  If $\beta_1=1$ and $\beta_2=\rho$, the oblique
            nullclines coincide and there is a continuum of equilibria rather
            than exactly one crossing.
          \end{enumerate}

          \item Direct differentiation gives
          \[
              \boxed{J(x,y)=
                  \begin{bmatrix}
                      1-2x-\beta_1y&-\beta_1x\\
                      -\beta_2y&\rho(1-2y)-\beta_2x
                  \end{bmatrix}}.
          \]

          \begin{enumerate}[label=\roman*.,leftmargin=2em]
            \item At $e_1=(1,0)$,
            \[
                J(e_1)=\begin{bmatrix}-1&-\beta_1\\0&\rho-\beta_2\end{bmatrix},
                \qquad \lambda_1=-1,\qquad \lambda_2=\rho-\beta_2.
            \]
            Thus $e_1$ is a stable node if $\beta_2>\rho$ and a saddle if
            $\beta_2<\rho$.  Biologically, a rare coyote population cannot
            invade the fox-only state when $\beta_2>\rho$, but can invade when
            $\beta_2<\rho$.  When $\beta_2=\rho$, one eigenvalue is zero and
            linearization is inconclusive.  A center-manifold calculation gives
            $y'=\rho(\beta_1-1)y^2+O(y^3)$ in the biologically relevant side:
            $e_1$ is attracting from the interior if $\beta_1<1$, repelling in
            the center direction if $\beta_1>1$, and belongs to a line of
            equilibria if also $\beta_1=1$.

            \item At $e_2=(0,1)$,
            \[
                J(e_2)=\begin{bmatrix}1-\beta_1&0\\-\beta_2&-\rho\end{bmatrix},
                \qquad \lambda_1=1-\beta_1,\qquad \lambda_2=-\rho.
            \]
            Hence $e_2$ is a stable node if $\beta_1>1$ and a saddle if
            $\beta_1<1$.  A rare fox population cannot invade the coyote-only
            state when $\beta_1>1$, but can invade when $\beta_1<1$.  For
            $\beta_1=1$, linearization is inconclusive.  Near the center
            direction, $x'=(\beta_2/\rho-1)x^2+O(x^3)$, so the point attracts
            from the interior if $\beta_2<\rho$, repels in that direction if
            $\beta_2>\rho$, and lies on a line of equilibria when
            $\beta_2=\rho$ as well.

            \item At the interior equilibrium, the nullcline identities reduce
            the Jacobian to
            \[
                J(e_4)=\begin{bmatrix}
                           -x_*&-\beta_1x_*\\
                           -\beta_2y_*&-\rho y_*
                \end{bmatrix}.
            \]
            Therefore
            \[
                \operatorname{tr}J(e_4)=-(x_*+\rho y_*)<0,
                \qquad
                \det J(e_4)=x_*y_*(\rho-\beta_1\beta_2).
            \]
            If $\beta_1<1$ and $\beta_2<\rho$, the determinant is positive,
            so $e_4$ is asymptotically stable.  If $\beta_1>1$ and
            $\beta_2>\rho$, the determinant is negative, so $e_4$ is a saddle;
            its stable manifold separates initial conditions leading to fox
            exclusion from those leading to coyote exclusion.
          \end{enumerate}

          \item
          \begin{enumerate}[label=\roman*.,leftmargin=2em]
            \item The model assumes a closed, well-mixed, spatially homogeneous
            system; continuous deterministic populations; constant parameters;
            logistic growth for each species in isolation; instantaneous,
            pairwise, density-dependent competition; no age or sex structure;
            no migration, time delay, seasonality, disease, prey dynamics, or
            environmental noise; and fixed carrying capacities.

            \item In
            \[
                r_iN_i\left(1-\frac{N_i}{K_i}\right)
                =r_iN_i-\frac{r_i}{K_i}N_i^2,
            \]
            $r_iN_i$ is density-independent growth and
            $-(r_i/K_i)N_i^2$ is intraspecific competition.  The terms
            $-\alpha_1N_1N_2$ and $-\alpha_2N_1N_2$ are the effects of
            interspecific encounters on foxes and coyotes, respectively.

            \item Abundance data and approximate low-density growth rates may be
            comparatively accessible, while carrying capacities remain
            habitat- and season-dependent.  The interaction coefficients
            $\alpha_1$ and $\alpha_2$ are usually hardest to identify because
            they require separating interspecific effects from habitat,
            resources, predation, and intraspecific density dependence.  The
            dimensionless parameters inherit this uncertainty.
          \end{enumerate}

          \item
          \begin{enumerate}[label=\roman*.,leftmargin=2em]
            \item There is a modeling mismatch if the equations are interpreted
            literally as a spatial description: a well-mixed ODE has no
            territories and cannot represent avoidance.  It can still describe
            average abundance effects after spatial behavior has been folded
            into the $\alpha_i$.  Also, coyote predation is not the same as
            symmetric competition: a predator--prey formulation would typically
            reduce fox growth but increase coyote growth, for example replacing
            the second interaction term by $+\gamma N_1N_2$ together with an
            appropriate coyote mortality/resource model.  Territories can be
            represented using spatial patches, diffusion/advection, or
            density-dependent encounter rates.

            \item Replacing red foxes by bobcats changes the parameter values and
            weakens the justification for a simple well-mixed competition model.
            Bobcats are solitary, use very large home ranges, are almost entirely
            carnivorous, and show strong habitat and seasonal effects; coyote
            interactions may also be important.  A better model could include
            spatial habitat classes, seasonal carrying capacities, shared-prey
            dynamics, and asymmetric competition.  The algebraic stability
            analysis is unchanged only if one retains the same two-species
            competition equations.

            \item Replacing coyotes by gray foxes also preserves the algebra only
            under the same modeling assumptions.  Gray foxes use brushy and
            deciduous woodland more than the open habitats preferred by red
            foxes, and their opportunistic diet overlaps only partly with that of
            red foxes.  Habitat partitioning should therefore reduce effective
            encounter rates or require separate patches.  Because coyotes may
            prey on young gray foxes, any retained coyote effect would be
            asymmetric rather than pure competition.  Thus one should
            re-estimate $r_2,K_2,\alpha_1,\alpha_2$ and consider explicit habitat
            and predation terms.
          \end{enumerate}
          \end{enumerate}
      \end{solution}
    \end{enumerate}



\end{document}